% s-014-p014-m002; original PDF 14, 15
\Aphorism{s-014-p014-m002}{3}
Das zweifellos Liebenswürdige an Menschen und Dingen ist ihre \Sperr{Idealität}, oder deutlicher, ihr Antheil an der Zukunft und ihre Loslösung von der Vergangenheit, ihre Duldsamkeit und Liebe zu noch unerprobten Möglichkeiten, ihre Freiheit von jeder Art „Ahnenstolz“ und Vergötterung des Gewesenen. „Vertrieben bin ich aus Vater- und Mutterländern; so liebe ich allein noch meiner Kinder Land, das unentdeckte, im fernsten Meere!“ Mich berauscht noch nach einem halben Jahrhundert die wundervoll schwärmerische Lenz- und Zukunftsstimmung, in der Richard Wagner in den Dresdener Maiaufstand hineingerieth — unschuldig, unbetheiligt, für seine eigentlichen Ziele ohne Hoffnung, ohne tieferen Zusammenhang mit dem „Volke“, nur von dem einen ächten Künstler-Instinct getragen: fort vom Vergangnen! hinaus in die Dämmerung des Kommenden! — Jenem seltsamen Manne, der als Mensch und Denker den schrittbreiten Pfad zwischen dem Erhabnen und dem Lächerlichen wandelt, fehlt es gerade an diesem \Sperr{Einen}, das eigentlich das Venerabile im Charakter ausmacht; er ist keiner von den Hoffenden, die irgend ein goldenes Kleinod von Wunsch und Sehnsucht dem Gange der Zeit vorauswerfen. Im Gegentheil hat \Sperr{Schopenhauer} eine gewisse Meisterschaft darin, der Zukunft ihren Sinn zu rauben und überall die Grundüberzeugung durchscheinen zu lassen, als ob mit seiner Philosophie die Welt enträthselt, ihr weiteres Fortbestehen also im höchsten Sinne überflüssig sei. Als philosophisches Dogma: weil die Menschheit bereits fähig ist, das Weltziel zu erreichen, d. h. den Willen zu verneinen, und \Foreign{natura nihil frustra facit}, darum muss die Menschheit die \Sperr{letzte} Objectivationsstufe des Willens sein. Welch ein tyrannischer Gedanke! welche Vergewaltigung der noch unerlebten Zukunft, welche eindeutige und ausschliessliche Interpretation des reichen vielgestaltigen Lebens! — Wenn nicht die Wahrheit selbst, so ist doch der Glaube an die gefundene Wahrheit in gefährlichem Masse lebensfeindlich und zukunftsmörderisch. Noch Keiner von denen, die sich mit Wahrheit begnadet wähnten, hat einen Augenblick gezögert, das grosse Finale oder den grossen Mittag oder irgend einen Endpunkt, Wendepunkt, Gipfelpunkt der Menschheit zu verkünden, d. h. jedesmal allem Künftigen sein Bild, seinen Stempel, seine Beschränktheit aufzuprägen.
\par

